\documentclass[11pt]{article}

\usepackage[margin=1in]{geometry}
\usepackage{booktabs}
\usepackage{hyperref}
\usepackage{url}
\urlstyle{tt}

% Prefer slightly looser interword spacing over text running past the
% margin; many lines here end in long unbreakable \texttt{} identifiers.
\setlength{\emergencystretch}{3em}

\title{11-768 HW2: Evaluating a Data-Visualization Agent}
\author{Zoe Xue \\ Andrew ID: zoex}
\date{}

\begin{document}
\maketitle

\paragraph*{Overview.} This report follows one design thread: \emph{every
decision that has a mechanical answer is taken away from the judge model
and made deterministic; the model is kept only for perception and
language understanding, and even there its claims are verified against
the pixels where physics allows.} Part 1 shows the released baseline is
uninformative (macro-MCC 0.027) and diagnoses two mechanisms. Part 2
fixes them --- a deterministic execution check and rubric-scoped,
structurally parsed questions --- reaching 0.448. Part 3 stress-tests
that validator with five adversarial self-authored tasks (0.105 there),
packages them as deterministic Harbor verifiers, and measures the exact
boundary where deterministic verification stops. Part 4 turns those
findings into a measured improvement campaign (frozen baselines, a
repeat-noise floor, pre-registered candidates, audits of every changed
prediction) whose product is the shipped validator:

\begin{center}
\begin{tabular}{lrr}
\toprule
validator & seed macro-MCC & authored macro-MCC \\
\midrule
released baseline & 0.0268 & --- \\
Part 2 & 0.4475 & 0.1047 \\
\textbf{shipped (Part 4)} & \textbf{0.5002} & \textbf{0.6210} \\
\bottomrule
\end{tabular}
\end{center}

Throughout, negative results are reported alongside positive ones: the
one family that trades score for evidence soundness, the components that
measurement rejected, and a reserved holdout on which the final
validator does \emph{not} beat its predecessor.

\section{Part 1: Examine Seed Agent Runs}

\subsection{Comparison of Ground-Truth and Validator Labels}
We ran the released baseline validator (\texttt{validator/baseline.py}) on all
110 seed runs and scored it against \texttt{seed\_labels.json}:

\begin{center}
\begin{tabular}{lrrrrr}
\toprule
Family & MCC & TP & TN & FP & FN \\
\midrule
\texttt{execution\_failure} & $-0.0022$ & 4 & 65 & 33 & 8 \\
\texttt{wrong\_data} & 0.0000 & 0 & 73 & 0 & 25 \\
\texttt{wrong\_chart} & 0.1096 & 3 & 27 & 0 & 68 \\
\texttt{hard\_to\_read} & 0.0000 & 0 & 59 & 0 & 39 \\
\midrule
macro-MCC & \textbf{0.0268} & & & & \\
\bottomrule
\end{tabular}
\end{center}

The baseline carries almost no signal, and its disagreements with the human
labels concentrate in two places. On \texttt{execution\_failure} it
\emph{over}-predicts: 33 of its 37 positive calls are false, while it misses
8 of the 12 genuinely failed runs. On the three content families it barely
predicts at all --- and an important control-flow fact shapes that number:
the baseline returns immediately after a positive execution verdict, so its
37 execution-positive runs never reach the content questions at all. Only
73 runs were actually asked the three content questions (219 questions, not
$3\times110$), and across those the baseline produced five raw content
positives (runs 059, 087, 096, 173, and 208, all \texttt{wrong\_chart}), of
which three are scored --- 096 and 208 are excluded because their ground
truth is a terminal \texttt{execution\_failure}. Against 135 positive
ground-truth content labels, \texttt{wrong\_data} and
\texttt{hard\_to\_read} are constant-NO (MCC 0 by definition) and
\texttt{wrong\_chart} keeps only the sliver of credit its three scored hits
earn.

\subsection{Validator Error Type 1: verdict-polarity confusion on the
execution question}
The baseline asks the negatively-framed compound question ``Did the agent
crash, exhaust its turns, or give up without producing a valid figure?'' and
takes any reply beginning with ``YES'' as an error. On 33 runs that produced
a perfectly valid figure --- among them \textbf{004, 022, 023, 034, 046, 062}
--- the recorded evidence \emph{contradicts the recorded verdict}: the reply
begins ``Yes'' and then continues, e.g.\ ``The agent successfully completed
the task. It created a Python script that generated a 2x2 grid of subplots
\ldots'' (run 004) or ``The agent successfully created the requested Sankey
diagram\ldots'' (run 046). The model is answering the affirmative question it
expected (``did the agent produce the figure?'') rather than the negative one
it was asked, and the prefix-based parser turns that polarity slip into a
false error. The mirror image appears on the 8 missed true failures
(\textbf{010, 053, 065, 089, 096, 168, 181, 208}): each of these run
directories contains no \texttt{figure.png} at all, so the label is
mechanical, yet none of them is predicted as an execution failure. Six have
entirely empty predictions (e.g.\ 010, 053), and two (096, 208) instead
carry a \texttt{wrong\_chart} error --- the judge went on to critique the
design of a figure that does not exist. Because the baseline does not
record negative replies, we cannot distinguish a genuine NO verdict from a
reply whose phrasing simply failed the leading-``YES'' parse; either way
the family is decided by fragile surface form rather than by the artifact.

\subsection{Validator Error Type 2: suppressed positives on the content
families}
For \texttt{wrong\_data}, \texttt{wrong\_chart}, and \texttt{hard\_to\_read},
the baseline asks one maximally broad question per family (``Does the figure
fail to follow any part of the requested chart design?'') with no definition
of what counts, and again keeps only replies that begin with ``YES''. The
result is near-total silence: 0/25, 3/71, and 0/39 scored true positives
recovered. The misses are not subtle judgment calls; they include defects
that are plainly observable in the figure and labeled with concrete
evidence, on runs that \emph{were} asked the content questions. In run
\textbf{009}, the x tick labels were requested rotated 30\,degrees, are
drawn horizontal, and overlap into unreadability --- \texttt{wrong\_chart}
and \texttt{hard\_to\_read}, neither flagged. In run \textbf{038}, the
`Non-smoker' series is drawn in the background colour and is invisible; in
run \textbf{016} the requested $<$0.5\% filter is ignored and the requested
legend is missing; in run \textbf{030} the title reads ``Shared axes''
instead of the requested exact text; in run \textbf{213} the Men's column
is plotted where the Women's column was requested. None of it is flagged.
Three mechanisms combine here. The question gives the judge no rubric for
what the family covers, so borderline YES cases default to a safe NO; the
prefix parse discards any positive expressed in a different surface form;
and error type 1 compounds the silence structurally --- the 37 runs the
baseline calls execution failures (run 004 among them, with its linear axis
where log-base-3 was requested and its omitted cubing step) never reach the
content questions at all, so a third of the dataset cannot produce a
content positive no matter what the figure shows.

\subsection{Proposed Validator Changes}
\textbf{For error type 1:} take the model out of the loop entirely.
\texttt{execution\_failure} is decidable from the artifact: a run failed
exactly when it left no decodable, non-empty \texttt{figure.png}. On the
seed set this deterministic rule agrees with the ground truth on all 110
runs (12/12 positives, 98/98 negatives) and cannot suffer polarity or
format errors.

\textbf{For error type 2:} replace each broad question with the assignment's
own rubric for that family (the full sub-criteria list), instruct the judge
to enumerate the task's explicit requirements and check them one by one
against the figure and the plotting code, and require a fixed two-line
answer tail (\texttt{VERDICT: YES/NO}, \texttt{EVIDENCE: \ldots}) that is
parsed structurally instead of by literal prefix, with a re-ask when the
reply truncates before its verdict.

\section{Part 2: Improve Your Validator}

\subsection{Implemented Changes}
Both changes are implemented in \texttt{validator/solution.py}.

\textbf{Change 1: deterministic execution check.} The baseline asks the VLM
whether the agent ``crashed, exhausted its turns, or gave up''.
\texttt{execution\_failure}, however, is defined by the artifact, not the
narrative: a run either left a decodable, non-empty \texttt{figure.png} or it
did not. The improved validator decides this family without the model: missing,
undecodable, or zero-sized figures are execution failures (evidence cites the
final trajectory message), and everything else is not. On the seed set this
rule reproduces the ground truth exactly (12/12 positives, 98/98 negatives).
It is also model-agnostic by construction, which matters because the grading
deployment uses the same fixed judge but a private set of runs.

\textbf{Change 2: rubric-scoped questions with structured verdicts.} The
baseline asks one broad yes/no question per family (e.g.\ ``does the figure
fail to follow \emph{any} part of the requested chart design?'') and accepts
the answer only if the reply literally starts with ``YES''. The improved
validator sends, for each remaining family, the assignment's own sub-criteria
for that family, an explicit instruction to check only what the task
\emph{requested} and to answer NO when the instructions left a choice open,
and a checklist procedure (enumerate the explicit requirements, then compare
each against the figure and the plotting code). Replies must end with
\texttt{VERDICT: YES/NO} and \texttt{EVIDENCE: ...} lines, which are parsed
structurally (with a robust fallback, and a re-ask when the reply was
truncated before reaching a verdict). The model sees the instructions
verbatim, the input files, a digest of the trajectory that preserves every
executed command (the plotting code) while clipping prose and tool output,
and the final figure. The baseline's reactive context-compaction and image
downscaling are retained for oversized runs.

\subsection{Seed-Set Performance vs.\ Baseline}
\begin{center}
\begin{tabular}{lrrrrrr}
\toprule
& \multicolumn{2}{c}{MCC} & \multicolumn{4}{c}{Improved (counts)} \\
Family & Baseline & Improved & TP & TN & FP & FN \\
\midrule
\texttt{execution\_failure} & $-0.0022$ & \textbf{1.0000} & 12 & 98 & 0 & 0 \\
\texttt{wrong\_data} & 0.0000 & 0.3372 & 17 & 51 & 22 & 8 \\
\texttt{wrong\_chart} & 0.1096 & 0.1387 & 20 & 23 & 4 & 51 \\
\texttt{hard\_to\_read} & 0.0000 & 0.3142 & 13 & 54 & 5 & 26 \\
\midrule
macro-MCC & 0.0268 & \textbf{0.4475} & & & & \\
\bottomrule
\end{tabular}
\end{center}

\subsection{Qualitative Analysis of Targeted Error Types}
\textbf{Error type 1 (execution polarity) is eliminated.} The deterministic
check scores a perfect 1.0: all 33 baseline false positives (e.g.\ 004, 022,
046) and all 8 false negatives (e.g.\ 010, 053) are gone, and no new errors
appear, because the decision no longer passes through the model at all.

\textbf{Error type 2 (suppressed positives) is strongly mitigated but not
solved.} Positive predictions on the content families went from 5 raw / 3
scored to 81, recovering 17/25 \texttt{wrong\_data}, 20/71
\texttt{wrong\_chart}, and 13/39 \texttt{hard\_to\_read} true positives;
every content-family MCC improved. Two causes contribute, and they should
not be conflated: the rubric prompts and structural parsing raise the yield
of the questions themselves, but change 1 also \emph{ungates} the content
questions --- of the 37 runs the baseline short-circuited as execution
failures, the 33 with valid figures are now all asked (the other four are
genuine failures and remain terminal), so part of the recall gain is
recovered coverage rather than better prompting. Where the questions fire, the
evidence is genuinely grounded: on run \textbf{013} the prediction reads
``uses the rainbow colormap instead of the required reversed plasma
colormap'' against the human label ``the instructions forbid
jet/rainbow/hsv, but the figure uses rainbow'' (a scored
\texttt{wrong\_chart} true positive), and on run \textbf{016} it reads
``includes versions with shares below 0.5\% (e.g., Chrome 69.0 at 0.43, IE
9.0 at 0.44)'' against the label's identical filter violation (a scored
\texttt{wrong\_data} true positive). Not every positive is this clean ---
runs 196 and 216 show the judge asserting defects (swapped group colours, a
stacked-vs-grouped confusion) that the figures do not support, and both are
scored false positives; hallucinated perceptual claims remain the main
precision cost, taken up in Parts 3 and 4.

\textbf{A representative remaining false positive.} Run \textbf{013} is
predicted \texttt{wrong\_data} with the evidence ``the task requires the
reversed plasma colormap, but the figure uses the rainbow colormap\ldots
explicitly prohibited''. The defect it describes is real, but a colormap
choice is chart \emph{style} --- the ground truth files it (and we define
it) under \texttt{wrong\_chart} --- so the same discovery is scored as a
\texttt{wrong\_data} false positive. Runs 030 and 051 show the same
signature (title text and colormap violations reported as data errors). The
judge finds the right defect and attributes it to the wrong family: the
rubric's exclusion lists (``do NOT report style here'') are not strong
enough to stop a salient violation from leaking into whichever question
notices it. A second false-positive mechanism appears in run \textbf{009}
(\texttt{wrong\_data}), where the model asserts that plotted matplotlib
version numbers are wrong and versions are omitted; the human label records
no data error, so this is a misreading of a dense timeline image ---
unverified perceptual claims stated as fact.

\textbf{A representative remaining false negative.} Run \textbf{030} is a
\texttt{wrong\_chart} false negative: the label's evidence is that the
figure's title reads ``Shared axes'' instead of the requested exact text
``Sharing x per column, y per row''. The validator \emph{did} report exactly
this defect --- but inside its \texttt{wrong\_data} answer, so
\texttt{wrong\_chart} recorded a miss and \texttt{wrong\_data} a false
alarm. One family confusion thus costs two categories at once. More broadly,
\texttt{wrong\_chart} keeps 51 false negatives against only 27 negative
runs; with truth this positive-heavy, the instruction to answer YES only
when confident suppresses exactly the recall this family needs. What is
missing from the validator is a cross-family arbitration step: verdicts are
made independently per family, so nothing re-files a discovered defect into
the family our definitions assign it to.

\section{Part 3: Design New Evaluation Tasks}

\subsection{Anticipated Generalization Weaknesses}
Reading \texttt{validator/solution.py}, the two families still delegated to
the VLM with the least corroborating signal are \texttt{wrong\_data} and
\texttt{hard\_to\_read}, which motivates two task classes.

\textbf{Class 1: derived-value aggregation} (\texttt{task\_class} =
``derived-value aggregation''). The validator judges plotted values from the
rendered image plus the agent's code. When correctness hinges on a multi-step
\emph{derived} quantity --- filter rows, then aggregate, then normalize or
weight, then order by the result --- the image alone cannot reveal a wrong
number (a bar of height 79 looks as plausible as one of height 49), so the
validator must effectively re-derive the answer from the code and input data.
That is exactly the kind of arithmetic-over-context that a 30B judge is weak
at, so near-miss aggregations (dropped filter, unweighted mean, raw counts
instead of shares, order computed from the wrong quantity) may slip through.

\textbf{Class 2: readability stress} (\texttt{task\_class} = ``readability
stress''). \texttt{hard\_to\_read} is judged purely by VLM perception of the
rendered image. Subtle legibility defects --- tick labels that collide or get
clipped by the image edge, a legend sitting on top of a data series,
neighbouring series in nearly identical colours --- are known VLM blind
spots, in both directions: a dense-but-legible chart invites false positives
while genuinely overlapping text can be glossed over. Tasks that force many
long tick labels or crowded legends probe this boundary directly.

\subsection{Task Design}
Five tasks (three in class 1, two in class 2); all data is synthetic and
small enough to hand-verify, and every instruction pins the title, axis
labels, ordering, and layout so correctness is human-reviewable.

\begin{itemize}
\item \textbf{\texttt{fine-totals}} (class 1; input \texttt{fines.csv}).
  Sum fines per library branch \emph{excluding waived rows}, bars in
  descending order of total. The data is adversarial: with the waived filter
  the totals are $\{$Eastside 53.5, Northgate 49.0, Harborview 42.5, Central
  37.5$\}$; without it they become $\{$Northgate 79.0, Central 77.5, Eastside
  55.75, Harborview 42.5$\}$ --- three of the four totals and the required
  bar order change (Harborview alone keeps 42.5, having no waived rows),
  so a dropped filter is visible in both data and order.
\item \textbf{\texttt{sensor-share}} (class 1; input \texttt{readings.csv}).
  Aggregate repeated per-site readings and plot each site's \emph{percentage
  share} of the grand total (exactly 17.5/21.0/31.5/30.0, summing to 100) as
  a horizontal bar chart, alphabetical top-to-bottom, with formatted value
  labels. Traps: plotting raw sums instead of shares, and matplotlib's
  \texttt{barh} default of drawing the first category at the \emph{bottom}.
\item \textbf{\texttt{weighted-latency}} (class 1; data inline).
  Request-weighted mean latency per service (90.0, 80.0, 56.0, and
  130.0~ms). The
  unweighted means (100/80/120/100) differ on three of four services, and the
  ``search'' service is a control where both agree.
\item \textbf{\texttt{stage-attendance}} (class 2; data inline). Twelve bars
  with long stage names as tick labels, which must be rotated 45°,
  right-anchored, and fully visible on a canvas of at least $14 \times 6$
  inches. Small default figures, missing rotation, or missing layout calls
  produce exactly the clipped/overlapping-label defects the class targets.
\item \textbf{\texttt{city-temps}} (class 2; data inline). Six crossing
  temperature lines with pinned per-city colours and a legend that must sit
  \emph{outside} the axes to the right without covering any line. Weak agents
  reliably leave the legend inside, where it occludes data --- our very first
  generated run (\texttt{city-temps\_\_qwen}) did exactly that.
\end{itemize}

\subsection{Agent Runs and Labeling}
Each of the five tasks was run once with each fixed agent
(Qwen2.5-Coder-3B-Instruct, Ministral-3-14B-Instruct-2512, GLM-4.7-Flash)
through the provided Modal agent runner: 15 runs, all under \texttt{runs/}.
Every run was reviewed against the assignment's evidence standard by reading
the instructions, the trajectory (including the plotting code), and the
rendered figure, and each label in \texttt{labels.json} cites its concrete
evidence.

The outcome splits cleanly by agent strength. Three Qwen-3B runs (on
\texttt{fine-totals}, \texttt{sensor-share}, and
\texttt{weighted-latency}) produced no figure at all and ended with
\texttt{LimitsExceeded}, each stuck in its own failure loop: the first
repeatedly invoked \texttt{sudo} in a container that has none (``sudo: not
found'', returncode 127); the second died on a Python
\texttt{SyntaxError}, a compound \texttt{for} statement placed after a
semicolon in a one-liner; the third alternated between a Python
\texttt{SyntaxError} and shell quoting errors (\texttt{/bin/sh}
complaining of an unexpected parenthesis). That is three
\texttt{execution\_failure} labels, each citing its own trajectory. Qwen's
other two runs produced figures with design errors: on
\texttt{stage-attendance} it drew the correct bars but left numeric
position ticks (0, 2, 4, \ldots, 10) instead of the required stage names
(\texttt{wrong\_chart}), and on \texttt{city-temps} it placed the legend
inside the axes against an explicit outside-right requirement
(\texttt{wrong\_chart}), where it covers the Farrow and Dunmore lines'
January--February points (\texttt{hard\_to\_read}). The mid-size agents
were stronger but fell into the designed ordering trap on
\texttt{sensor-share}: both GLM and Ministral computed the percentage
shares correctly and then let \texttt{barh} draw Alder at the bottom where
the task demands it at the top (\texttt{wrong\_data}, wrong order).
Ministral's run also lacks the required ``Site'' y-axis label in the saved
image --- its script does call \texttt{set\_ylabel}, but the label lands
outside the saved region, and a required text element entirely outside the
image counts as missing (\texttt{wrong\_chart}). The remaining eight runs
satisfied every requirement and are labeled acceptable.

Two labeling decisions deserve comment. First, the inverted
\texttt{sensor-share} bar order is filed under \texttt{wrong\_data} (not
\texttt{wrong\_chart}) because the family definitions place ``wrong order:
\ldots sequences appear in a different order from the one requested'' under
data. Second, for \texttt{city-temps\_\_qwen} the same legend produces two
distinct errors --- its placement violates an explicit design requirement,
and, independently, it occludes data points --- which is exactly the
design-versus-legibility distinction the two families encode; the occlusion
was confirmed by inspecting a crop of the legend region.

\subsection{Validator Performance on New Tasks}
The Part 2 validator, run unchanged on the 15 authored runs
(the pre-iteration authored predictions file in \texttt{output/}):

\begin{center}
\begin{tabular}{lrrrrrr}
\toprule
Family & MCC & TP & TN & FP & FN & support ($+$/$-$) \\
\midrule
\texttt{execution\_failure} & 1.0000 & 3 & 12 & 0 & 0 & 3/12 \\
\texttt{wrong\_data} & $-0.3162$ & 0 & 6 & 4 & 2 & 2/10 \\
\texttt{wrong\_chart} & $-0.1741$ & 0 & 8 & 1 & 3 & 3/9 \\
\texttt{hard\_to\_read} & $-0.0909$ & 0 & 10 & 1 & 1 & 1/11 \\
\midrule
macro-MCC & \textbf{0.1047} & & & & & \\
\bottomrule
\end{tabular}
\end{center}

With supports this small each error moves MCC a long way, but the failures
are systematic, not noise, and land exactly in the two designed weakness
classes. Three patterns account for every mistake. \textbf{(1) Hallucinated
colour violations:} on all three \texttt{city-temps} runs the validator
reported that the lines are ``blue instead of tab:blue, orange instead of
tab:orange'' --- but \texttt{tab:blue} \emph{is} matplotlib's default blue,
so the claim describes no visible difference. It poisoned the two clean runs
(false \texttt{wrong\_data} on GLM, false \texttt{wrong\_data} +
\texttt{wrong\_chart} on Ministral) and displaced the real errors on the
Qwen run, which was predicted \texttt{wrong\_data} instead of its true
\texttt{wrong\_chart} + \texttt{hard\_to\_read}. \textbf{(2) Missed order
inversion:} on both \texttt{sensor-share} runs the alphabetical-from-top
requirement is violated by \texttt{barh}'s bottom-up default, and the
validator saw nothing wrong --- the class-1 hypothesis (weak spatial/derived
reasoning) confirmed on ordering rather than on values, which the agents
computed correctly. \textbf{(3) Cross-family mis-filing with a perceptual
hallucination:} on \texttt{stage-attendance\_\_qwen} it correctly noticed
the tick labels are not the stage names but filed the defect under
\texttt{wrong\_data}, and added a \texttt{hard\_to\_read} claim that ``the
full stage names are cut off'' --- there are no stage names in that figure
at all. The agents' failures (three qwen \texttt{LimitsExceeded} runs, the
legend and tick-label mistakes, the \texttt{barh} inversion) are weaknesses
of the \emph{agents}; every scoring mistake above is a weakness of the
\emph{validator}, and the three patterns directly motivate Part 4.

\subsection{Harbor Packaging}
Each of the five tasks is packaged under \texttt{harbor/tasks/<task\_id>/}.
The verifiers use three layers: S1-style artifact checks (a decodable,
non-blank PNG of sane size; the input CSV unmodified, held to a sha256
literal; a plotting script left behind), an exact comparison of
\texttt{plotted\_values.json} against expected values derived by hand and
written into the test as literals, and one re-execution of the delivered
script in a scrubbed copy of \texttt{/app} (delivered artifacts are
snapshotted first), which must exit 0 and reproduce a sidecar identical to
the delivered one. All five tasks score 1.0 with \texttt{-a oracle}. The
worked example was also run as required: oracle 1.0, \texttt{nop} 0.0, and
two of its mutants (\texttt{sum\_crates} 0.0; \texttt{illegible} 1.0).

\textbf{Mutant scoring 0: \texttt{fine-totals/mutants/no\_filter}.} A
plausible dropped-clause error: it sums every row, ignoring the waived
filter. Reward 0. It passes every artifact and re-execution check and is
rejected by exactly one assertion,
\texttt{test\_sidecar\_declares\_}\allowbreak\texttt{exactly\_the\_expected\_values},
which
reports the declared $\{$Northgate 79.0, Central 77.5, Eastside 55.75,
Harborview 42.5$\}$ against the hand-derived key.

\textbf{Mutant scoring 1.0: \texttt{stage-attendance/mutants/illegible}.}
Every number is correct, the sidecar matches the key, and --- deliberately
--- every structurally checkable clause of the instructions is satisfied:
the canvas is the required $14\times6$ inches, the tick labels are the full
stage names, rotated 45\,degrees and right-anchored, and all text lies
inside the saved image. What is wrong is purely perceptual: every text
element is set at 2\,pt, roughly three pixels tall at 100\,dpi, unreadable
by any person. Reward 1.0. Our verifier accepts it because its checks read
file metadata, JSON values, and re-execution behaviour, never the rendered
pixels; the S1 gates (image $>100$\,px on a side, $>1000$ bytes, $>4$
colours) are all cleared. We state the boundary carefully: this acceptance
partly reflects checks we did not write --- a verifier that instrumented
matplotlib's artist tree, as \texttt{harbor/example} does, could read the
font sizes off the figure and reject this particular mutant with a
threshold. But legibility in general (labels that collide under rotation,
low-contrast text, a legend occluding data) has no per-clause deterministic
test, which is why the assignment's own accepted example works the same
way. To make the contrast measured rather than asserted, we ran our
validators on this mutant's figure (with its instructions and a minimal
trajectory), twice each. The Part 2 validator flags
\texttt{hard\_to\_read} --- the right family --- but with a fabricated
mechanism (it calls the tick labels ``clipped'' when they are merely far
too small). The final validator never accepts the mutant either, but on
both repeats it files the problem elsewhere: a \texttt{wrong\_data} claim
fabricated from the axis text it cannot read, and a \texttt{wrong\_chart}
answer that in one repeat describes the true defect exactly (``font size
2 \ldots\ the text is unreadable'') while filing it in the wrong family.
Its edge-verification stage, which requires a clipping claim to leave ink
at the named image edge, correctly refuses the fabricated ``clipped''
mechanism --- and with it, the lucky-family credit the Part 2 validator
earned. Text that is present, unclipped, and simply too small sits in the
residual blind spot of the deterministic verifier and both VLM validators
at once; we report that plainly rather than claiming the boundary
crossed.

\textbf{Could a deterministic per-task verifier decide each family?}
\texttt{execution\_failure}: yes --- file existence and PNG decodability are
fully deterministic (our validator already decides it this way).
\texttt{wrong\_data}: for an \emph{authored} task, largely yes --- the sidecar
plus hand-derived literals decides the plotted values, and re-execution ties
them to the script; what it cannot see is a figure whose drawn geometry
disagrees with the sidecar, which needs either savefig instrumentation (as in
\texttt{harbor/example}) or perception. \texttt{wrong\_chart}: only
partially --- structural facts (orientation, axis limits, tick text) are
recoverable with matplotlib instrumentation, but only per-task, per-clause;
each explicit design requirement needs its own assertion.
\texttt{hard\_to\_read}: mostly no --- individual proxies exist (a
font-size threshold read off an instrumented artist tree would catch our
own accepted mutant), but legibility as a property (labels colliding under
rotation, contrast against a background, a legend occluding data) lives in
the rendered pixels, admits no per-clause test, and is where the accepted
mutant sits relative to the value-level checks we wrote.

\textbf{Why a per-task verifier cannot be the grading mechanism.} A Harbor
verifier encodes the answer key for one known task: hand-derived literals, a
schema for the sidecar, sha256 digests of the inputs, assertions matched to
each clause of one instruction. The private grading set contains tasks we
never saw, so there are no literals to write, no schema to pin, and no
clause-by-clause assertions to prepare. Grading there requires a single
function that generalizes across arbitrary instructions --- deriving the
expected chart from the task text, data, trajectory, and image at run time
--- which is precisely the validator, not a fixed test suite.

\section{Part 4: Iterate on Your Design}

\subsection{Additional Improvements Motivated by Part 3}

\paragraph{Why a campaign rather than one change.} Part 3 exposed three
failure patterns: colour hallucinations, missed order inversions, and
cross-family mis-filing. Our first response was a composite change ---
deterministic colour filters plus a model arbitration pass --- but a
subsequent review found both components defective: the colour filter
deleted valid findings (a genuine ``the requested tab:blue series is
red'' loses its sentence to word co-occurrence; ``the title is missing
and the line colour is incorrect'' loses the title defect too), and the
arbitration parser could silently erase findings for families a
truncated reply never answered. Patching by intuition had produced
plausible-looking components that measurement rejected, so we re-ran
Part 4 as a measured campaign instead.

\paragraph{Method.} We froze the incumbent (predictions, scores, source
hashes, a manifest of run/input hashes), added optional per-call
diagnostics (disabled by default; reads no labels; changes no
decisions), built a fixed 39-run diagnostic subset enriched with hard
cases from both released and authored runs, and measured the incumbent
twice on it and twice on the authored set, giving a repeat-noise floor
of about $\pm$0.02--0.03 macro-MCC that every promotion is judged
against. Each candidate was pre-registered (hypothesis and plausible
regression written before running), gated on local regression tests,
evaluated on the subset, and audited finding by finding; the full ledger
ships in \texttt{experiments/LEDGER.md}. The diagnostics answered the
first question directly: of $\sim$38 content-family misses on the
subset, 34--35 were never detected by the first-stage questions ---
recall was bounded by \emph{detection}, not by filtering or arbitration
--- which set the candidate order below.

\paragraph{Candidates and decisions.} (Subset macro-MCC; incumbent
repeats 0.268/0.289.)

\begin{center}
\begin{tabular}{p{8.2cm}rl}
\toprule
candidate (one component each) & subset & decision \\
\midrule
remove destructive colour filter; arbiter may not invent verdicts & 0.269 & adopt (correctness) \\
checklist \texttt{wrong\_chart} prompting & 0.370 & adopt \\
readability prompt alone & 0.379 & adopt only with edge check \\
deterministic mapping replaces model arbiter & 0.432 & adopt \\
composed + edge verification + fact layer & \textbf{0.610} & champion \\
\quad + dual-filing dedup (post-champion) & replay & adopt \\
\quad + code-literal contrast check (post-champion) & audit & adopt \\
\bottomrule
\end{tabular}
\end{center}

\paragraph{Adopted changes, in pipeline order.}
\begin{enumerate}
\item \textbf{Correctness repairs.} The colour filter is disconnected
  (four of five required-preservation regression cases fail against it),
  and the arbiter can no longer emit a verdict for a family no reviewer
  proposed with fabricated ``confirmed by arbitration'' evidence. Subset
  score unchanged within noise --- kept because verified correctness
  repairs are not traded for score.
\item \textbf{Checklist \texttt{wrong\_chart} prompting.} The broad
  ``check each requirement'' instruction becomes an enumerated checklist
  (at most 12 items: quote, expected, observed, source image/code,
  status pass/fail/unknown) with an explicit image-over-code rule --- a
  property set in code but absent from the saved image is absent. Subset
  \texttt{wrong\_chart} $-0.045 \to +0.279$ (TP 5$\to$15), recovering
  legend, axis-label, orientation, axis-limit, and colorbar defects,
  including \texttt{sensor-share\_\_ministral}'s off-canvas label.
\item \textbf{Readability prompting with deterministic edge
  verification.} \texttt{hard\_to\_read} no longer answers under the
  ``explicit requirement'' confidence gate (legibility defects are
  rarely explicit clauses); it walks element groups with CHECK lines and
  must name the element, its location, and a visible mechanism. Alone
  this recovered real occlusion/overlap/invisibility defects but
  amplified ``partially cut off'' fabrications (11 new false alarms); an
  edge-ink audit showed only one clipping claim in the whole subset run
  was physically real. A deterministic verifier therefore strikes any
  clipping claim whose named image edge holds no ink. Replayed on the
  same run: \texttt{hard\_to\_read} $+0.107 \to +0.444$, false positives
  16$\to$1.
\item \textbf{Deterministic family mapping replaces model arbitration.}
  An ablation replaying \emph{identical} recorded findings under three
  configurations scored: direct $+0.378$; deterministically re-filed
  $+0.436$; model-arbitrated $+0.370$. The arbiter erased more true
  positives than it repaired, so findings are re-filed by rule (drawn
  order $\to$ \texttt{wrong\_data}; missing/wrong title, label, legend,
  colormap, scale $\to$ \texttt{wrong\_chart}; partial clipping $\to$
  \texttt{hard\_to\_read}) and the arbitration call is gone --- one
  model call per positive run saved.
\item \textbf{Deterministic fact layer.} An explicit DPI request is
  compared with the PNG's own metadata (fires on exactly the three true
  violations across all 125 development runs; silent on the twelve
  compliant DPI-mentioning runs), and literal number series in the
  instructions are compared with literal lists in the last successful
  figure-writing command (fires on exactly one run, 077's re-ordered
  series; sorting requests suppress the order rule). Both audited
  offline to zero false positives before inclusion.
\item \textbf{Dual-filing dedup.} One real defect reported by two family
  questions at once costs a guaranteed false positive in one of them
  (``sd instead of sem'' filed as both \texttt{wrong\_data} and
  \texttt{wrong\_chart}). Near-duplicate cross-family errors --- a
  similarity gate calibrated on 75 real evidence pairs, with genuinely
  distinct coexisting defects sitting safely below it --- collapse toward
  the family their combined language names.
\item \textbf{Code-literal contrast check.} Several
  \texttt{hard\_to\_read} truths share one mechanism: an element drawn
  in a colour too close to the background, with the colour sitting in the
  agent's code as a literal. Colour kwargs from every successful command
  (excluding edge colours, legend-proxy handles, unlabeled reference
  lines, translucent fills, masking patches, and dark-canvas runs) are
  compared against the rendered corner background; a WCAG-style contrast
  ratio below 2.5:1 files \texttt{hard\_to\_read}. After three audit
  iterations it fires on 8 runs across all 125 development runs --- every
  one a \texttt{hard\_to\_read} truth --- with zero false positives.
  Replaying both of these last two components over the three fixed seed
  sweeps of the pre-revision champion raised macro-MCC by a consistent
  $+0.045$ to $+0.055$ (about twice the repeat-noise floor), with
  \texttt{hard\_to\_read} moving from $0.08$ to $0.26$--$0.29$.
\end{enumerate}

\subsection{Final Evaluation on Seed and New Tasks}
The designated final sweeps of the shipped validator, against the Part 2
version's designated sweeps:

\begin{center}
\begin{tabular}{lrrrr}
\toprule
& \multicolumn{2}{c}{Seed runs} & \multicolumn{2}{c}{Authored runs} \\
Family & Part 2 & Final & Part 2 & Final \\
\midrule
\texttt{execution\_failure} & 1.0000 & 1.0000 & 1.0000 & 1.0000 \\
\texttt{wrong\_data} & 0.3372 & \textbf{0.3553} & $-0.3162$ & \textbf{0.0756} \\
\texttt{wrong\_chart} & 0.1387 & \textbf{0.3061} & $-0.1741$ & \textbf{0.4082} \\
\texttt{hard\_to\_read} & 0.3142 & \textbf{0.3395} & $-0.0909$ & \textbf{1.0000} \\
\midrule
macro-MCC & 0.4475 & \textbf{0.5002} & 0.1047 & \textbf{0.6210} \\
\bottomrule
\end{tabular}
\end{center}

Seed \texttt{wrong\_chart} recall rose from 23/71 to 50/71 true
positives, and \texttt{hard\_to\_read} --- the family Part 4's first
revision traded away --- recovered to 0.3395 (14/39 true positives, 5
false positives) once the code-literal contrast check joined the
deterministic layer. Repeat variation is reported in full rather than summarized
away: before selection, finalist and incumbent were each run twice more,
fresh and interleaved under matched serving:

\begin{center}
\begin{tabular}{lrrrr}
\toprule
macro-MCC & \multicolumn{2}{c}{seed} & \multicolumn{2}{c}{authored} \\
 & r1 & r2 & r1 & r2 \\
\midrule
Part 2 incumbent & 0.4389 & 0.4700 & 0.3552 & 0.2808 \\
pre-revision champion & \textbf{0.4800} & \textbf{0.4939} & \textbf{0.6787} & \textbf{0.6529} \\
shipped (designated / repeat) & 0.5002 & 0.5093 & 0.6210 & 0.6416 \\
\bottomrule
\end{tabular}
\end{center}

On seed, the champion's worse repeat beats the incumbent's better one.
The two final deterministic components (dedup, contrast) were
additionally measured by exact replay over the three fixed pre-revision
seed sweeps --- $0.4800/0.4931/0.4939 \to 0.5338/0.5367/0.5411$, a
consistent $+0.045$ to $+0.055$ --- which isolates their effect from the
fresh-sweep variance of the model-judged families.

Two results are reported against ourselves. First, the
\texttt{hard\_to\_read} story is a deliberate two-step trade. The
edge-verification stage initially dropped the family to 0.08: we
investigated every lost true positive across the comparison sweeps, and
all twelve were fabricated clipping claims that happened to land on runs
whose real defect is a different mechanism (overlapping tick labels, an
invisible series, upside-down text, a legend covering tick labels) ---
unsound credit, removed on purpose. The code-literal contrast check then
rebuilt recall on verifiable evidence (008's \texttt{\#f4f4f4} labels,
200's white scatter series), landing at 0.3395 --- above the Part 2
version's 0.3142, now with every clipping or contrast claim physically
checked. The remaining 25 misses are non-clipping mechanisms the judge
still overlooks (e.g.\ 009's overlapping tick labels). Second, before
tuning we built and reserved a fresh 26-run validation suite (12
correct/faulty pairs plus two execution failures; constructed and
visually verified by us, not human-reviewed course labels; provenance
and scores in \texttt{experiments/}) and evaluated it exactly once
after selecting the champion: incumbent macro $+0.7248$, champion
$+0.7108$ --- statistically a tie inside the noise floor, with the
champion perfect on readability (4/4, no false positives) and more
false-positive-prone on clean charts elsewhere. The dedup component was
subsequently designed partly from that suite's observed errors, so it is
now development data and \emph{no untouched holdout remains}; the
development-set gains do not demonstrate private-set superiority, and
the defensible claim stays: comparable-or-better classification with
substantially sounder evidence, at one model call less per positive run.
Our engineering stretch target of 0.60 seed macro-MCC was not reached
(designated 0.5002; replay band 0.534--0.541); the subset result
(0.610) shows the pipeline can get there on hard cases, and the open
\texttt{hard\_to\_read} and \texttt{wrong\_chart} recall gaps are
the identified remainder.

\subsection{Qualitative Analysis of Self-Authored Tasks}
On the designated authored sweep the shipped validator scores 0.6210,
catching \texttt{city-temps\_\_qwen} exactly (\texttt{wrong\_chart}
+ \texttt{hard\_to\_read}, the legend-over-lines mechanism named
correctly), \texttt{sensor-share\_\_ministral}'s order inversion as
\texttt{wrong\_data}, and \texttt{stage-attendance\_\_qwen}'s
missing stage names; the twin inversion on
\texttt{sensor-share\_\_glm} was caught in the previous designated
sweep but missed in this one --- the run-to-run instability of that
detection, noted below, cuts both ways.

\textbf{A false positive.} \texttt{city-temps\_\_ministral} (a correct
run) is again flagged, this sweep as \texttt{wrong\_data}: ``Brinport
is shown in orange instead of tab:orange (which is correct), but the
main [violation]\ldots'' --- evidence that concedes its own claim. The
trajectory contains \texttt{'Brinport': 'tab:orange',} and the drawn
lines match it; the claim asserts a difference no pixels can show.
This is the known cost of removing the colour filter: that filter was
measurably destructive (it also deleted genuine colour defects), so we
took the principled trade and this hallucination channel stays open. The
second residual FP channel is re-derivation: \texttt{fine-totals\_\_ministral}
is accused of breaking the descending order its own code
(\texttt{sorted(totals, key=..., reverse=True)}) demonstrably satisfies
--- the judge re-ranks the four close totals in its head and trusts its
arithmetic over the image.

\textbf{A false negative.} \texttt{sensor-share\_\_ministral}'s
\texttt{wrong\_chart} label is still missed. The script calls
\texttt{ax.set\_ylabel('Site')}, but the label lies entirely outside the
saved image; a judge that weighs the code sees the requirement satisfied.
The final validator can catch this mechanism class --- its checklist's
image-over-code rule recovered this very run on the diagnostic subset ---
but detection is unstable run to run, and the designated sweep missed it
(that sweep ran without diagnostics, so we cannot say which checklist
line failed). Noticing an \emph{absence} in the image against a
\emph{presence} in the code remains the hardest case for the judge.

\subsection{Further Changes}
Everything beyond the adopted changes is infrastructure, shipped for
reproducibility rather than score: the optional diagnostics layer
(\texttt{HW2\_DIAG\_DIR}), the frozen Part 2 validator
(\texttt{validator/solution\_part2.py}; the pre-iteration numbers
regenerate with \texttt{--solution validator.solution\_part2}), the
experiment ledger and per-candidate frozen sources under
\texttt{experiments/} (e00 incumbent freeze, candidate shas, the
arbitration ablation, the reserved-suite builder and its labels), and
regression tests covering the filtering and malformed-reply cases that
started the campaign. Cost: the shipped validator makes 3 judge calls
per valid-figure run (plus at most one bounded truncation re-ask each;
no arbitration call), and 0 calls for missing-figure runs; the full
110-run seed sweep completes in roughly 35--45 minutes on the course
deployment. Known open limitations, measured and not patched around:
residual dual-filing (the dedup stage collapses the pattern only above
its calibrated similarity gate; a sub-gate pair survives in the
holdout's scatter-styled run), re-derivation and misbinding
hallucinations on correct charts (now the dominant false-positive
channel --- e.g.\ the fine-totals order re-derivations in the
designated sweep), the non-clipping \texttt{hard\_to\_read} recall
gap (25 seed misses, e.g.\ 009's overlapping tick labels), and
tiny-but-unclipped text, which evades the deterministic verifier and
both VLM validators alike. The natural next candidate --- targeted crop
verification for proposed legibility defects --- is specified in the
ledger but was not pursued after the campaign's stopping review.

\subsection{Reproducibility}
\sloppy
Every number in this report regenerates from the submission. The Part 1
table: run \texttt{validator.baseline} through \texttt{validator.runner}
on \texttt{artifacts/} and score against \texttt{seed\_labels.json}. The
Part 2 tables: the frozen module \path{validator/solution_part2.py},
selected with the runner's \texttt{--solution} flag; its designated
predictions are \path{output/pre_iteration_*.json}. Parts 3--4: the
shipped \path{validator/solution.py}; designated predictions are
\path{output/final_seed_predictions.json} and
\path{output/authored_predictions.json}, scored with
\texttt{workflow score} (against \path{labels.json} for authored
runs). Harbor rewards: \texttt{workflow check-submission} re-runs every
oracle and both mutants. The experiment ledger --- per-candidate
hypotheses, frozen sources with SHA-256s, subset scores, the arbitration
ablation, and the reserved-suite builder and labels --- is under
\path{experiments/}; local gates run offline via
\path{experiments/regression_tests.py} and
\path{experiments/arbitration_ablation.py}. Judged sweeps are
nondeterministic at the margin (batched serving), so comparisons should
use the reported repeat bands rather than single sweeps.
\par\fussy

\end{document}
